\providecommand{\C}{\mathbb{C}}
\providecommand{\Z}{\mathbb{Z}}
\providecommand{\Q}{\mathbb{Q}}
\providecommand{\Sp}{\mathrm{Sp}}
\providecommand{\SO}{\mathrm{SO}}
\providecommand{\Auts}{\mathrm{Aut}_{s}}
\providecommand{\Fix}{\mathrm{Fix}}
\providecommand{\fg}{\mathfrak{g}}
\providecommand{\Bor}{\mathrm{Bor}}
\providecommand{\Grit}{\mathrm{Grit}}
\providecommand{\kch}{\kappa_{\mathrm{ch}}}
\providecommand{\kBKM}{\kappa_{\mathrm{BKM}}}
\providecommand{\kcat}{\kappa_{\mathrm{cat}}}
\providecommand{\kfib}{\kappa_{\mathrm{fiber}}}
\providecommand{\phiZO}{\phi_{0,1}}
\providecommand{\GammaSix}{\Gamma_{6}^{+}}
\providecommand{\DeltaOneSix}{\Delta_{1}^{(6)}}
\providecommand{\gSix}{\fg_{\DeltaOneSix}}

\section*{$N=6$ primitive CHL denominator $\DeltaOneSix$}
\label{sec:N6-primitive-chl-denominator}

\paragraph{Geometric datum.}
Let $g_6\in\Auts(K3)$ be a symplectic automorphism of order $6$, Mathieu
class $6A$, cycle shape $1^2 2^2 3^2 6^2$. Nikulin's fixed-point counts
are
\[
  \#\Fix(g_6)=2,\qquad
  \#\Fix(g_6^2)=6,\qquad
  \#\Fix(g_6^3)=8 .
\]
The topological Lefschetz formula gives traces on $H^2(K3)$
\[
  \operatorname{tr}(g_6\mid H^2)=0,\qquad
  \operatorname{tr}(g_6^2\mid H^2)=4,\qquad
  \operatorname{tr}(g_6^3\mid H^2)=6 .
\]
Writing
\[
  H^2(K3,\C)\cong
  \C^{a}\oplus\C_{-1}^{b}
  \oplus(\C_{\zeta_3}\oplus\C_{\zeta_3^{-1}})^c
  \oplus(\C_{\zeta_6}\oplus\C_{\zeta_6^{-1}})^d
\]
gives
\[
  a+b+2c+2d=22,\quad a-b-c+d=0,\quad
  a+b-c-d=4,\quad a-b+2c-2d=6 .
\]
Hence $(a,b,c,d)=(6,4,4,2)$. Therefore
\[
  r_6=\mathrm{rk}\,H^2(K3,\Z)^{g_6}=6,
  \qquad
  \mathrm{rk}\,H^2(K3,\Z)_{g_6}=16 .
\]
The invariant rank on full Mukai cohomology is $8$.
The CHL/umbral Niemeier anchor is
\[
  N_6=D_4^{\oplus6},
\]
root system $6D_4$, Coxeter number $6$, with umbral group
$G^{(6D_4)}=3.\mathrm{Sym}_6$ of order $2160$.

The quotient $K3/\langle g_6\rangle$ therefore has two stabiliser-$\Z/6$
orbits, two stabiliser-$\Z/3$ orbits, and two stabiliser-$\Z/2$ orbits:
\[
  K3/\langle g_6\rangle
  \quad\hbox{has Du-Val type}\quad
  2A_5+2A_2+2A_1 .
\]
The Burnside quotient Euler characteristic is
\[
  e(K3/\Z_6)=\frac{24+2+6+8+6+2}{6}=8,
\]
and the exceptional curves contribute $16$, giving Euler characteristic
$24$ after crepant resolution. The Niemeier lattice $A_5^4D_4$ is the
neighbouring Coxeter-number-$6$ entry with outer group of order $144$;
the $6A$ umbral character selects $6D_4$.

Let $t_6\in E[6]$ have exact order $6$ and set
\[
  X_6=(K3\times E)/\langle(g_6,t_6)\rangle .
\]
For $1\leq a\leq5$ the translation $t_6^a$ acts freely on $E$.
Hence $(g_6^a,t_6^a)$ acts freely on $K3\times E$, and $X_6$ is a
smooth CHL threefold with trivial canonical bundle. This is the
$K3\times E$ Calabi-Yau convention used here; the strict
simply-connected convention is a different geometric category. The
Du-Val calculation above belongs to the surface quotient
$K3/\langle g_6\rangle$; the threefold quotient is free.

\paragraph{Automorphic convention.}
Throughout this note $\DeltaOneSix$ denotes the primitive CHL $N=6$
Borcherds denominator. The orthogonal model uses
\[
  L_6^{\mathrm{aut}}=2U\oplus\langle -12\rangle
\]
of signature $(2,3)$; $\GammaSix$ is the corresponding extended
paramodular group, with character $v_6$. The divisor is the reflective
Heegner divisor cut out by the polar terms
$q^0\zeta^{\pm1}$ of the Jacobi seed below, equivalently the diagonal
Humbert divisor in the paramodular model, with multiplicity one in this
normalisation. This is the primitive CHL normalisation, distinct from
the Gritsenko-Cl\'ery diagonal-divisor atlas and from
Cl\'ery-Gritsenko-Hulek Niemeier square-root forms unless a row map is
specified.

\paragraph{Twined Jacobi seed.}
The primitive CHL convention fixes
\[
  (c_N(0))_{N=(1,2,3,4,6)}=(10,8,6,4,2),
  \qquad
  (\kBKM(\Phi_N))_{N=(1,2,3,4,6)}=(5,4,3,2,1).
\]
The $6A$-twined K3 weak Jacobi form of weight $0$ and index $1$ has
\[
  \phi^{(g_6)}_{0,1}(\tau,z)
  =(\zeta^{-1}+2+\zeta)+q(\cdots),
  \qquad \zeta=e^{2\pi iz}.
\]
In the primitive BKM-denominator normalisation,
\[
  c_6(0)=[q^0\zeta^0]\,\phi^{(g_6)}_{0,1}=2 .
\]
The same $q^0$ head gives
\[
  [q^0\zeta^{-1}]\,\phi^{(g_6)}_{0,1}
  =[q^0\zeta]\,\phi^{(g_6)}_{0,1}=1,
\]
which fixes the divisor multiplicity above.
The character $v_6$ on $\GammaSix$ has order $6$ and is determined by
the $6A$ multiplier of the twined Jacobi form.

\paragraph{Borcherds lift to paramodular weight $1$.}
The primitive CHL denominator table contains the chiral-half forms
\[
  \Delta_5,\quad
  \Delta_{4}^{(2)},\quad \Delta_{3}^{(3)},\quad
  \Delta_{2}^{(4)},\quad \Delta_{1}^{(6)}
\]
of weights $5,4,3,2,1$. Their scalar dyonic normalisations are separate
square branches. Borcherds' singular-theta weight theorem gives
\(\mathrm{wt}(\Bor(\phi))=c(0)/2\) for the chosen input normalisation.
Therefore
\[
  \DeltaOneSix=\Bor\bigl(\phi^{(g_6)}_{0,1}\bigr)
  \in S_1(\GammaSix,v_6),
  \qquad
  \mathrm{wt}\,\DeltaOneSix=\frac{c_6(0)}{2}=1 .
\]
Thus
\[
  \kBKM(\DeltaOneSix)=1 .
\]
The Borcherds product has the form
\[
  \DeltaOneSix(Z)
  =e^{-2\pi i(\rho_6,Z)}
  \prod_{\alpha\in\Delta_+}
  \bigl(1-e^{-2\pi i(\alpha,Z)}\bigr)^{c^{(g_6)}(\alpha)},
\]
where the exponents are the Fourier coefficients of
$\phi^{(g_6)}_{0,1}=\sum c^{(g_6)}(n,\ell)q^n\zeta^\ell$ under the
standard identification with $(L_6^{\mathrm{aut}})^\vee$, and $\rho_6$
is the leading-exponent vector of the product. The product is an
automorphic target identity; source realisation requires separate Hall
and orientation data.

\paragraph{Recognition data.}
The Borcherds product supplies automorphic target coefficients. A
level-$6$ denominator presentation requires a recognition datum
\[
  \mathscr R_6=(M_6,\iota_6,\Delta^{\mathrm{re}}_6,W_6,\rho_6,
  \Delta^{\mathrm{im}}_6,m_6,\epsilon_6)
\]
consisting of a Lorentzian lattice, a real-simple-root chamber, a Weyl
group, an embedding $\iota_6\colon M_6\hookrightarrow
(L_6^{\mathrm{aut}})^\vee$, imaginary simple roots with signed
multiplicities, and a parity rule. Let $\mathcal S_6$ be the finite
multisets of pairwise orthogonal imaginary simple roots allowed by
$m_6$, including the empty multiset, put $\alpha_S=\sum_{a\in S}a$, and
normalise $\epsilon_6(\varnothing)=1$. The bracket, parity blocks, and
Weyl alternating sum live in this recognition datum.

\paragraph{Denominator identity target.}
Given $\mathscr R_6$, the Weyl-Kac-Borcherds denominator identity for
the corresponding target $\gSix$ has product side
\[
  \DeltaOneSix(Z)
  =e^{-2\pi i(\rho_6,Z)}
  \prod_{\alpha\in\Delta_+(\gSix)}
  \bigl(1-e^{-2\pi i(\alpha,Z)}\bigr)^{
    \operatorname{smult}_{\gSix}(\alpha)}
\]
\[
  =\sum_{w\in W_6}\det(w)
  \sum_{S\in\mathcal S_6}
  \epsilon_6(S)\,
  e^{-2\pi i(w(\rho_6+\alpha_S),Z)} .
\]
For $\alpha=(n,\ell,m)$ in the positive chamber,
\[
  \operatorname{smult}_{\gSix}(\alpha)=c^{(g_6)}(nm,\ell),
\]
after transport by $\iota_6$. The polar coefficients of the Jacobi seed
determine the imaginary simple-root packet, while all Fourier
coefficients determine the root supermultiplicities. The
character $v_6$ and the parity rule $\epsilon_6$ fix the signs in the
superdenominator. The product identity is target-side data; a compact
BPS Lie algebra theorem also requires the Hall source, CoHA
multiplication, orientation package, and bracket comparison.

\paragraph{Subscripted invariants.}
The finite etale cover $K3\times E\to X_6$ and the Todd-class formula
give
\[
  \kcat(X_6)=\frac{1}{6}\chi(\mathcal{O}_{K3\times E})
  =\frac{1}{6}\chi(\mathcal{O}_{K3})\chi(\mathcal{O}_E)
  =\frac{1}{6}(2\cdot0)=0,
\]
equivalently $\chi(\mathcal O_{K3\times E})=6\chi(\mathcal O_{X_6})=0$.
The Hodge supertrace is also immediate. Since $g_6$ is symplectic and
$t_6$ is a translation,
\[
  H^{0,\bullet}(X_6)=
  \bigl(H^{0,\bullet}(K3\times E)\bigr)^{\langle(g_6,t_6)\rangle}
  \cong
  \C\oplus\C[-1]\oplus\C[-2]\oplus\C[-3],
\]
and therefore
\[
  \kch(X_6)=\sum_q(-1)^qh^{0,q}(X_6)=0,
  \qquad
  \kBKM(\DeltaOneSix)=c_6(0)/2=1 .
\]
The equivariant K3 fixed-lattice rank is
\[
  r_6=\mathrm{rk}\,H^2(K3,\Z)^{g_6}=6 .
\]
It is a level-$6$ lattice datum. The untwisted $K3\times E$ fibre
invariant remains
\[
  \kfib(K3\times E)=\mathrm{rk}\,\widetilde H(K3,\Z)=24.
\]
Thus the primitive CHL row contributes
\[
  (\kcat(X_6),\kch(X_6),\kBKM(\DeltaOneSix);r_6)=(0,0,1;6),
\]
while the untwisted compact $K3\times E$ invariants are
\[
  \kcat(K3\times E)=0,\qquad \kch(K3\times E)=0,\qquad
  \kch^{\mathrm{Heis}}(K3\times E)=3,\qquad
  \kBKM(\Delta_5)=5,\qquad \kfib(K3\times E)=24 .
\]

\paragraph{Compact-source comparison.}
A compact Hall theorem for $X_6$ starts with orientation data, HN
completion, and reduced equivariance. Its positive half has the form
\[
  Y^+_{\mathrm{Hall}}(X_6)=H^*_{\mathrm{eq}}
  \bigl(\mathcal{M}^+_{\mathrm{eff}}(X_6),\phi_W\bigr)
\]
for effective BPS moduli with vanishing-cycle sheaf $\phi_W$. Since
$X_6$ is a CY$_3$, the native $\Phi$-output is $E_1$-chiral; Drinfeld
double, vertex envelope, derived centre, and module-centre layers are
further operations. An $E_2$-structure belongs to those layers. The
native algebra $A=\Phi_3(\mathcal C_{X_6})$ remains $E_1$.

The local toric reference point is
\[
  \mathrm{CoHA}(\C^3)=Y^+ .
\]
The full $\mathcal W_{1+\infty}$ object is obtained only after
Drinfeld-double, centre, completion, or evaluation operations.

The automorphic denominator defines the target formal supercharacter
\[
  \operatorname{sch}\mathcal{F}_{\DeltaOneSix}(\tau,\sigma,z)
  =-\frac{1}{\DeltaOneSix(Z)}.
\]
A scalar reduced curve-counting theorem for $X_6$ is a separate
statement. In the Bryan-Oberdieck and Oberdieck CHL normalisation, its
scalar denominator is the square
\[
  \Phi^{\mathrm{scal}}_6(Z)=(\DeltaOneSix(Z))^2 .
\]
The scalar prefactor is fixed by the Bryan-Oberdieck chamber, boundary
terms, and variable convention. Such a scalar identity gives a protected
index; the Hall bracket, orientation character, and source-side
denominator algebra are additional structures. A source theorem
identifying the protected character of $Y^+_{\mathrm{Hall}}(X_6)$ with
$\mathcal{F}_{\DeltaOneSix}$ requires these data. A bracket-level
isomorphism
\[
  \fg_{\mathrm{BPS}}^{(6)}\cong\gSix
\]
requires the Davison-Meinhardt integration map, parity-refined PBW
comparison, and compatibility with Kontsevich-Soibelman wall crossing
in the level-$6$ chamber. This BKM/Hall-Drinfeld branch is separate from
the K3 Mukai self-mirror Yangian branch.

\paragraph{Normalisations.}
The Niemeier label is fixed by the umbral character:
$6D_4$ pairs with $\chi_{6D_4}\colon3.\mathrm{Sym}_6\to\C^\times$ and
the shadow of $H^{(6)}$. The lattice $A_5^4D_4$ has the same Coxeter
number and outer group of order $144$, giving a different umbral
assignment.

The Gritsenko-Cl\'ery diagonal-divisor catalogue is the separate
triple-indexed atlas
\[
 (t,N;k)\in\{(1,1;5),(2,1;2),(1,2;3),(3,1;1),(1,3;2),
 (4,1;\tfrac12),(1,4;\tfrac32),(2,2;1)\}.
\]
Its twice-weight tuple is
\[
  (10,4,6,2,4,1,3,2).
\]
The primitive CHL ladder used here has tuple
\[
  (c_N(0))_{N=(1,2,3,4,6)}=(10,8,6,4,2).
\]
The common entry of the two systems is $\Delta_5$ at level $1$.
In the Gritsenko-Cl\'ery atlas, weight one occurs at $(3,1;1)$ and
$(2,2;1)$; the primitive CHL orbifold order $6$ is a different index.

The weight is fixed by the Borcherds input constant:
\[
  k_6=\mathrm{wt}\,\DeltaOneSix=c_6(0)/2=1.
\]
Formulae from other Niemeier substitutions compute different
automorphic normalisations.

The level is paramodular:
\[
  \GammaSix\subset\Sp_4(\Q),
\]
with character $v_6$ of order $6$. The Siegel congruence subgroup
$\Gamma_0^{(2)}(6)\subset\Sp_4(\Z)$ is a different modular surface.

\paragraph{The $N=1$ anchor.}
The primitive Igusa normalisation is
\[
  D_5=64^{-1}\Delta_5,
  \qquad
  \Phi_{10}=\Delta_5^2,
  \qquad
  \kBKM(\Delta_5)=10/2=5.
\]
For the untwisted $K3\times E$ problem,
\[
  \kcat(K3\times E)=0,\qquad
  \kch(K3\times E)=0,\qquad
  \kch^{\mathrm{Heis}}(K3\times E)=3,\qquad
  \kBKM(\Delta_5)=5,\qquad
  \kfib(K3\times E)=24,
\]
with each value attached to its own construction.

References: Borcherds 1998 \emph{Invent.\ Math.}~$132$, Theorem~13.3;
Gritsenko 1999 \emph{St.\ Petersburg Math.\ J.}~$11$, Theorem~1.2;
Gritsenko-Nikulin 1998 Part II, Theorem~2.1;
Govindarajan-Krishna 2010 \emph{JHEP}~$05$, Table~1;
Gritsenko-Cl\'ery 2008, Theorems~1.4 and~3.1;
Mukai 1988 \emph{Invent.\ Math.}~$94$; Kondo 1998
\emph{Duke Math.\ J.}~$92$; Nikulin 1980 Theorem~4.15;
Cheng-Duncan-Harvey 2014 \texttt{arXiv:1307.5793};
Oberdieck 2018 \texttt{arXiv:1706.09377};
Bryan-Oberdieck 2018 \texttt{arXiv:1811.06102}.
